\chapter{Arquitectura y Diseño del Pipeline de Super--Resolución}
\label{ch:arquitectura}

Este capítulo describe la arquitectura del \textbf{Subsystema de Super--Resolución
basada en Inteligencia Artificial (SAI)}, detallando el flujo de datos, los
componentes funcionales y las decisiones de diseño adoptadas para la aplicación
de técnicas de super--resolución como postprocesamiento de imágenes adquiridas
por el CMOS Microscopy Payload (CMOSM).

La arquitectura se diseña con un enfoque modular, permitiendo la evaluación
experimental de distintos modelos de super--resolución sin afectar la
trazabilidad ni la integridad de los datos originales.

% -------------------------------------------------------------------
\section{Visión general del pipeline}

El pipeline del subsistema SAI se estructura como una cadena de procesamiento
secuencial, en la cual la imagen capturada por el CMOSM constituye la entrada
principal. El procesamiento se realiza fuera del flujo de adquisición y no
interfiere con la operación nominal del instrumento.

De manera general, el flujo de datos comprende las siguientes etapas:

\begin{enumerate}
    \item Ingesta de imagen y metadatos.
    \item Preprocesamiento y normalización.
    \item Inferencia mediante modelo de super--resolución.
    \item Postprocesamiento y generación de salida.
    \item Registro de trazabilidad.
\end{enumerate}

% -------------------------------------------------------------------
\section{Flujo de datos}

\subsection{Entrada del sistema}

La entrada al subsistema SAI corresponde a imágenes adquiridas por el CMOSM,
almacenadas en formatos sin compresión (PNG o TIFF) y acompañadas de metadatos
experimentales.

Los datos de entrada incluyen:
\begin{itemize}
    \item Imagen base calibrada (dark frame y flat field aplicados).
    \item Resolución nativa del sensor.
    \item Parámetros de adquisición (exposición, ganancia, iluminación).
\end{itemize}

La imagen de entrada se preserva de manera íntegra durante todo el proceso.

% -------------------------------------------------------------------
\subsection{Preprocesamiento}

Antes de la inferencia, las imágenes son sometidas a una etapa de
preprocesamiento, cuyo objetivo es asegurar compatibilidad con los modelos de
super--resolución evaluados.

Las operaciones de preprocesamiento incluyen:
\begin{itemize}
    \item Conversión a escala de grises, cuando aplica.
    \item Normalización de intensidad a un rango definido.
    \item Recorte o padding para ajuste a tamaños de entrada del modelo.
\end{itemize}

Estas operaciones no modifican la información estructural de la imagen, sino que
facilitan la inferencia del modelo.

% -------------------------------------------------------------------
\section{Modelos de super--resolución evaluados}

El subsistema SAI evalúa modelos de super--resolución basados en aprendizaje
profundo, seleccionados por su disponibilidad, madurez y desempeño demostrado
en tareas de realce visual.

\subsection{Modelos tipo ESRGAN}

Se consideran variantes de \textit{Enhanced Super--Resolution Generative
Adversarial Networks (ESRGAN)} como modelos principales de evaluación. Estos
modelos se caracterizan por:

\begin{itemize}
    \item Arquitectura generativa orientada a la mejora perceptual.
    \item Capacidad para realzar bordes y texturas.
    \item Uso extendido en tareas de super--resolución visual.
\end{itemize}

Los modelos ESRGAN se emplean exclusivamente con fines de realce visual, sin
atribuirles validez metrológica.

\subsection{Modelos alternativos}

De manera complementaria, se consideran arquitecturas convolucionales más
simples, como variantes de U--Net o redes residuales, con el objetivo de comparar
resultados y evaluar compromisos entre complejidad y estabilidad visual.

% -------------------------------------------------------------------
\section{Inferencia}

La inferencia del modelo constituye la etapa central del pipeline. En esta fase,
la imagen preprocesada es alimentada al modelo de super--resolución seleccionado,
generando una imagen de mayor resolución aparente.

El proceso de inferencia se caracteriza por:
\begin{itemize}
    \item Ejecución fuera de línea (off--line).
    \item Tiempo de procesamiento dependiente de la plataforma utilizada.
    \item Ausencia de retroalimentación al sistema de adquisición.
\end{itemize}

No se contempla inferencia en tiempo real ni a bordo del satélite en la fase
actual del proyecto.

% -------------------------------------------------------------------
\section{Postprocesamiento}

Posterior a la inferencia, la imagen generada puede someterse a ajustes menores
de visualización, tales como:
\begin{itemize}
    \item Reescalado de intensidad para visualización.
    \item Conversión de formato de salida.
\end{itemize}

En todos los casos, la imagen procesada se almacena como un producto derivado,
manteniendo siempre disponible la imagen original de entrada.

% -------------------------------------------------------------------
\section{Trazabilidad y gestión de datos}

El subsistema SAI implementa mecanismos de trazabilidad para asegurar la relación
entre imagen original y resultado procesado. Cada imagen procesada se asocia a:

\begin{itemize}
    \item Identificador de la imagen de entrada.
    \item Modelo de super--resolución utilizado.
    \item Parámetros relevantes del procesamiento.
\end{itemize}

Esta información se registra en archivos de metadatos independientes, permitiendo
auditoría y reproducibilidad del procesamiento.

% -------------------------------------------------------------------
\section{Consideraciones de diseño}

Las decisiones de diseño del pipeline SAI se fundamentan en los siguientes
criterios:

\begin{itemize}
    \item Separación estricta entre adquisición y procesamiento.
    \item Preservación de la imagen original sin modificaciones.
    \item Modularidad para evaluar distintos modelos de IA.
    \item Enfoque conservador respecto a la interpretación de resultados.
\end{itemize}

Este enfoque permite aprovechar técnicas modernas de aprendizaje profundo sin
comprometer la integridad científica del sistema de adquisición.

% -------------------------------------------------------------------
\section{Resumen de la arquitectura}

La arquitectura del subsistema SAI proporciona un marco estructurado y trazable
para la aplicación de técnicas de super--resolución como herramienta de apoyo a
la interpretación visual. El diseño adoptado es coherente con el rol secundario
del SAI dentro del proyecto INTISAT y sienta las bases para futuras extensiones
del pipeline.
